\chapter{附录}
.， 附 … 录 I 张量(包括向量)运算基础 1.正交直线坐标基矢量的运算公式 (1)克罗内克尔 (Kronecker) 符号 S 任意两个正交坐标基矢量点积用 (Jij 表示，称为 kronecker (J: 仇 i = 10 ${\rm •}$ i 手 j'; = ei ${\rm •}$ e; = ( 11. i = j 式中 i .j 是自由指标，此可写作 z

\begin{gather*}
(Jll = (J22 = (J33 = 1\\
i.j = 1.2.3\\
(J23 = (J32 = (J12 = (J21 = (J13 = (J31 = 0
\end{gather*}

(2) 置换符号 任意两个正交单位向量叉积可表示为 ei X ej = E 训 el 式中 E 圳称为置换符号，也称为里奇 (Ricc i)符号，其数值如下: 41 ( 1. 1) ( 1 .2) ( 1.3) 式( 1 . 1)对应 i.j .k 中有 2 个或 3 个相同指标 z 式<1.2) 对应 i.j .k 按 123123 … 顺序排列 F 式( 1.3) 对应 i.j .k 按 132132 …逆序排列.由此可知 Eijl 中任意两个自由

附录 指标对换，相应的量相差一个负号，如 E'23 -一句 '3. 2. 坐标基的转换公式 两直角坐标基 ei 和衔'间有以下关系 z

\[e , = \alpha jie j' e i = \alpha ijej\]

直角坐标基变换系数的是对应坐标轴间的方向余弦(见表 1) ，并有以下性质 z

\[\alpha .J=4 ,  \alpha m\alpha  jm = \delta '1 , \]

襄 1 坐标轴闰方向余弦 t. t 2

\[, t. \alpha  . . \alpha .2\]

e2 , a2. a22

\[t3 \alpha 31 a32\]

3. 坐标变换公式 任意两个坐标基 e ， 和 e/I中的位置向量 x 和 x' 有以下的转换公式 z X ; = 吨iXj' I ， = $\alpha$ ijX j 4. 向量和张量的定义 根据标量、向量和张量在坐标变换中的性质来定义向量和张量. (1)标量 t ,

\[\alpha 13\]

a:!'3

\[\alpha 33\]

和空间点有关的物理量 $\psi$( 町，巧 ， X 3 ) ，若它在直角坐标变换.r; =$\alpha$Jerf 时满足 WI1$\cdot $r2 $\cdot $h)=WZI' ， z; ， rifh 即其值不因坐标的刚性转动而变，则伊为一个标量. (2) 向量 向，句，向是直角坐标系 OX ， XZ X3 中的 3 个量，如果它们根据叫=的 aj 变换到另 一坐标系 OX.'X:X; 中的 3 个量叫， a; ，叫，则此 3 个量定义 一 个向量 4. (3) 张量 $\rho$阳是直角坐标系 OXIXZ X3 中的 9 个量，如果它按照下列公式

\[iJ=\alpha i/ Qj",  \theta 'm\]

转换为另一个直角坐标系 OX.'X:X; 中的 9 个量抖，则此 9 个量称为笛卡儿二阶张 量，简称二阶张量.推广到 n 阶张量定义如下. 设在每一个坐标系内给出 3" 个量 $\rho$扎1 ， 1旷2 叮.

\[\rho ' 1 ;2" ' ; , ,  -\alpha ')11\alpha ' 2 1 2 \cdots  \alpha  i"i"P jli 2\]

转换，则此 3" 个量定义了一个 n 阶张量.由此可看出 :n=O 时，张量的分量只有一 个，且满足 : p'= $\rho$ 的关系，因此是一个标量，故可视标量为零阶张量.当 n=l 时，张

392 附录 量有 3 个分量，且满足 zA=$\alpha$'1 1n] $\rho$m ，的关系，因此它是一个向量，可视向量为一阶 张量. 5. 张量识别定理 定理1. 1 若 $\rho$风，川1 1 2" 飞叫叩J${\rm 々}$2 "'J与m 和\{任圭意 m 阶张量 q飞3川"句2 "'J与翩的内积 $\rho$; ， ;Z""R;'li 2 "'\} m. qj 1i2… im = 1;1 性… 恒为 n 阶张量，则 $\rho$'1 '2 …'" I\} I 1'1. "'， 'm 必为 m 十 n 阶张量，上式中重复下标求和. 定理1. 2 若 $\rho$'1 '2 """和任意 m 阶张量 Qhj2000jm 的乘积 p勺 i'!.... ，，， qjjJ 2 ...\}m - 1;l i"!."'i ,. ;JI J2"'\}m 恒为 n+m 阶张量，则矶 '2""" 必为 Tl 阶张最.

\subsection*{例 1.1  任何向量 ai 和克罗内克尔符号 a ij 乘积为 ai =éJij 酌，闲此根据张量识别}

定理知克罗内克尔符号岛是二阶张量.坐标基变换系数也是二阶张量，又称转动张 量，常用 Q 表示.转动张量是正交的，因为它和它的转置张量 QT 之和、是克罗内克尔 符号 ${\rm e}$Ji， . 同理，用张量识别定理可以证明 E i) k 是三阶张量. 定义1. 1 并矢 z 两向量 a=aiei 和 b=h氏的下列乘积运算定义为并矢，并用 ab 表示 z

\[ab = a ,h ,eie ,\]

用张量识别定理也可证明并矢是二阶张量. 6. 向量代数运算公式

\[a ::l::: b = aiei ::l::: bie , = (a , ::l::: h, )e,\]

a ${\rm •}$ b = aie 曹 ${\rm •}$ h,e, = aib ,e, ${\rm •}$ ej = a,hja ij = a ,b, e, el e3 a X b = aie; X hjej = aibjei X ej = $\alpha$ 川. = I a , a2 a3 bl fh b3 a. (b X c) = (a X b) ${\rm •}$ c = c. (a X b) = (c X a) ${\rm •}$ b a X (b X c) = (a ${\rm •}$ c)b - (a ${\rm •}$ 的 c (aXb). (cXd) = (a.c)(b.d) 一 (b.c)(a.d) 7. 二阶张量代数运算 (1)二阶张量代数运算公式(用并矢表示)

\[ab ::l::: cd = (a ;bj ::l::: cidj)eiej\]

a ${\rm •}$ bc = (a ${\rm •}$ b)c = c(b ${\rm •}$ a) ab ${\rm •}$ cd = a(b ${\rm •}$ c)d = (b ${\rm •}$ c)ad = ad(c ${\rm •}$ b)

附录 c. ab ${\rm •}$ d = (c ${\rm •}$ a)(b ${\rm •}$ d) = (b ${\rm •}$ d)(c ${\rm •}$ a) = (d ${\rm •}$ b)(a ${\rm •}$ c)

\[abXc=a(bXc)\]

(2) 二阶张量分量的变换式 按张量定义 z

\[B = bi;eie j = b;'j' e;, ej'\]

式中 b'J 和 11 ，/ 分别是对应于两个坐标系中的张量分量，其对应的坐标变换为 bi,j' e\textasciitilde{}. (b,.. e $\iota$，) ${\rm •}$ e; = btm(e:' et)(e> e..) , i' ,j' = 1. 2, 3 biJ = e, ${\rm •}$ (bt' m'e:e\textasciitilde{}) ${\rm •}$ ej = b(m'(e , ${\rm •}$ e;)(e J ${\rm •}$ ${\rm e}$m) , ;, j = 1.2.3 基向量点乘结果参见表1. 8. 向量场的微分运算公式 (\})哈密顿( Hamilton) 算子v 晗密顿算子是一个具有微分和向最双重运算的算子，哈密顿算子在直角坐标系 中的表达式为 J . . d . . r7 V=i 一十 j 一 +k 一 =e 一一r7.l" . " iJ y . \_. rJ z -, rJ.r, (2) 定义 ${\rm 1}$标量和向量的梯度 标量 $\psi$ 的梯度在直角坐标系中表示为 grad$\psi$=iif 十 j 尹 +k 学=吟

\[.  .1- "y "Z\]

向量 b 的梯度在直角坐标系巾表示为 ${\rm •}$ rJ b , ${\rm •}$ cJ b , ${\rm •}$ cJ b gradb = i 了+ j .\textasciitilde{}- + k .\textasciitilde{}- = Vb

\[rJ.l" "y <l Z\]

${\rm 2}$向量和二阶张量的散度 在直角坐标系中向量 a 在 M 点的散度用 diva 表示: diva =华+华+些- V. a- 些L rI .r rly rlZ rI.l", 在直角坐标系中二阶张量 B 的散度用 divB 表对 z r1 B. . r7 B.. r7 B divB =一\textasciitilde{}+一-.v +一-y = V. B rJ.1' . Jy . rJ z 向世的散度是标量.二阶张量的散度是向量. ${\rm 3}$向量的旋度 向世的旋度用 rota 或 curla 表示.在直角坐标系中可表示为 rota= i( 华 -EEi)+j( 丛一生二 )+k( 丛一旦旦) 飞rJ V ()z I 飞 ()Z d.1' I 飞 d .1' dV I

\[=VM=MhiEL cJ:r.\]

(3) 常用哈密顿算子运算微分公式

\begin{gather*}
V(\psi  + I\theta ) = VfJ +VI\theta \\
V(\psi  I\theta ) = fJ VI\theta  +l\theta V\psi 
\end{gather*}

V F( $\psi$) = F'( 仰伊，例如 Vf\{J (r) = $\psi$'$\omega$ 手 V. (a+ 的 =V.a+V.b V ${\rm •}$ (\textasciitilde{}口 a) = f\{JV , a + Vf\{J' a V ${\rm •}$ (a X b) = b ${\rm •}$ (V X a) - a ${\rm •}$ (V X b)

\begin{gather*}
V X (a + b) = V X a + V X b\\
V X (fJ a) = fJV X a + VfJ X a
\end{gather*}

V X (a X b) = (b ${\rm •}$ V )a 一 (a. V )b + aV ${\rm •}$ b - bV ${\rm •}$ a V (a ${\rm •}$ 的 = (b ${\rm •}$ V )a + (a ${\rm •}$ V )b + b X (V X a) + a X (V X b) la 尸(a-V)a=V-E一--aX (VXa)

\begin{gather*}
V. (VfJ) = \partial .fJ\\
V X (VfJ) = 0\\
V. (V X a) = 0
\end{gather*}

V X V X a = V (V ${\rm •}$ a) - ${\rm ð}$.a V. (f\{JV I${\rm þ}$) = $\psi$${\rm ð}$.1${\rm þ}$ +V伊. VI${\rm þ}$

\[\partial . (fJ I\theta ) = 1\theta \partial .fJ + fJ\partial . \psi + 2 VfJ' VI\theta \]

E 高斯公式和斯托克斯公式 1.广义高斯 (Gauss) 公式 附录 如果 A 是空间体积 r 的封闭曲面，物理量 a 或 v 在 r+A 上一阶偏导数连续，则 有以下体积分和面积分间的等式 z LV' 附=扣 ${\rm •}$ adA A LVf\{J dr = L时 M LVX 附= Ln XadA 以上三式称为广义高斯公式.式中 n 为微元曲面 dA 外法线方向单位向量(参见 附图 1). 上式还可用于 a 为二阶张量的情况.

附录 395 (a) (b) 附图 1 (a) 说明高斯公式用阁 I (b) 说明斯托克斯公式用图 式州、

\begin{gather*}
-my\\
--mHH
\end{gather*}

门 wA 分=积'町的。，'广 V 界盯川叩顿密哈町，" ill V 0 时 = ffa 0 ndA ill V X 时 = ffn X adA ill<tl O V) 阶 = ff<tI 0 n)adA - illav 0 时r

\[ill~\psi dr = ff::dA = ffn 0 VqJ dA\]

ill\textasciitilde{}附 = ff\textasciitilde{}:dA = ffn 0 Va ill < qJ $\tilde{t}$${\rm þ}$ + V qJ$\psi$o Vt${\rm þ}$驯)油d击r = ffqJ 兰dA ill (qJ $\tilde{t}$${\rm þ}$ 一向)d$\tau$ = ff( 忐一 d)" 1| 问 1 2 dr =恍dA (其中 $\psi$ 满足句= 0) 3. 斯托克斯 (Stokes) 公式 如 L 为曲面 A 的边界，且为可缩曲线，向量 a 在 A+L 上一阶偏导数连续，则有 以下回路线积分和面积分间的等式(参见附图1): Ln o 何$\times $山 =PLa 0 也 式中 n 为微元曲面dA的法线方向，其指向与积分周线"L 方向符合右手螺旋法则，如

396 附录 附图 1 所示. E 正交曲线坐标系 1.坐标基矢量 空间曲线坐标系由气组空间由面的交线组成，如果空间三组曲面的交线相互垂直， 则构成正交曲线坐标系(参见附图 2). 宅间曲面在直角坐标系 Cr.y.z) 中的方程写作: q, = q, (.r .y.z). q2 = q2 (.l" .y.Z). q.l = q" Cr.y.z) (皿.1)

\begin{gather*}
(X.".=) ,.\\
M (1'.1/1.=) q,
\end{gather*}

、 (R. fJ.申1 ,. Y (a) (b) 附图2 正交刷线坐标系 (al 直角坐标、柱坐标相球坐f;J、; (hl .般正交 rth 线坐伪、 正交曲线坐标系的摹矢量 e ， 由坐标面 q$\iota$i = q, (.r. 成'应用曲面的法向量公式，基矢量 Vqi (m.2) I Vq, I 基矢量还可以用另一种方法表示，由坐标面方程组(山.1)可导出直角坐标系和 宅间曲线坐标系之间的转换公式 z .r = ,r(q , .q, .q3)' y = y( 酌 .q ， .q:I)' Z = Z(CJ, .q, .q,,) (U1. 3) 在这组方程巾，令 q2.q :l 等于常数，就可以得到 q ， 坐标线的参数方程，而基矢量沿坐 标线的切线方向.由式( 1lI. 3). (.l' .y.z) 空间的点向量可表示为 R = .r (q , ,q, .q,,)i + y(q , ,(/, .q :l )j 十 ::(q ， ， 矶. (/:1 ) k ( III . 4) 曲线坐标 q ， 的单位切向蛙.即坐标基

\[e, =\]

2王 i+ 主卫 j+ 旦王 KJq , - . r7 q,oI . rJ q,

\[(25)2+(iff+(if)z\]

附录 可见只要给定曲线坐标系与直角坐标系的变化关系，就可计算出任意宅间点上的曲 线坐标系的基矢量 e ， : 与$\cdot $+izj+3三k

\[e, =\]

d qi Oqi "qi (考 f + (老 f + (考 f (皿， 5) 2 ，拉梅 (Lamc) 系数 正交曲线坐标线上的微元弧 l支出，与坐标量 dq ，的比值称为拉梅系数，用 h ， 表 I c1 R , 1... - $\tilde{L}$- ${\rm •}$ 示.坐标线的微元弧 K ds,= I 一句， I (对 F 标 z 不求和) .则拉梅系数为I rJ q, \_..' I dSi I  R I 11 rJ.r \textbackslash\{\} 2 , 1 习 v\textbackslash\{\}2 ， l rJ z\textbackslash\{\}2 h ， =7一=七一 I=A/\{'\textasciitilde{}\textasciitilde{}1 +\{'$\tilde{Jl}$ + 七\textasciitilde{} 1. i = 1. 2.3 (皿， 6) dqi I ('q , I ''V飞 oq ， I 飞 "qi I 飞 oq ， I 若已知 q ， =qi (.r .y.z). 可用 F 列方法求 h ， : hi - 坐!....=土=一止一 =\_1dqi dqi e,. V q, I V qi I d 五， 町-学习+组j+ 学$\tilde{k}$ rJ.I (J Y (IZ 3 ，基矢量的导数 由于 e ， 是 q ， 的函数.所以正交曲线坐标系中基矢量的导数不全等于零，这是它 和直角坐标系的主要差别 .jt 计算公式如下 z 1 rJ h 手土 =77句， ( i 手) .式Ij I i.j 不做求和运算)

\begin{gather*}
"(/, 11 , "q,\\
iJel ( 1 all I - I 1 Jh 1 -  
\end{gather*}

32: 一-\textbackslash\{\}瓦石72 l 瓦石71

\[c1 e., 1 1 rJ II ., , 1 ah.,  \]

5石=- \textasciitilde{}，，-;- c1 q;el T 瓦而 e :I )

\[iJe:I ( 1 ,111 :1 - I 1 c1h:, -  \]

JL= 一町石-1:" 2 ---r-瓦石75j 柱在坚标、球坐标系巾坐标基的导数公式可由以上公式计算得到. fJJJ ][. I 柱坐标系 柱坐标系的坐标挝 (r.O.z) 和直角坐标(.r l ..2" 2 ${\rm •}$ .Z':! )间的关系为 基矢量表示为 r = (.d 十 .ri\}4. 。 =arctan 主王 ， Z = .I:i .1' 1 (皿， 7a) (囚， 7b) (囚， 7c) (囚， 7d)

\[e. = X1 e, + X2 ,. - -C' I -,-- -1:" 2 , r r\]

拉梅系数表示为 九= 1, 基矢量的导数表示为 X2 \_ , X1 e8 =-$\tilde{e}$1 寸-\textasciitilde{}2' r r

\[h8 = r , h. = 1\]

Je, Je8 \_ \_ E百一时， J(\} -,

\[e. = e3\]

Je, \_ Je8 \_ Je. \_ Je. \_ Je, \_ Je8 \_ Je. \_ " Jr Jr Jr J(\} Jz Jz Jz 例:1 .2 球坐标系 球坐标系的坐标量 (R ，()， $\psi$) 和直角坐标 (X1 ,X2 ， X3) 间的关系为

\[x. .  =arccos h ,,\]

(xf + x\textasciitilde{} + x\textasciitilde{} )吉 R = (xf + x\textasciitilde{} + x\textasciitilde{} 汁， (/J = arctan -=二 X1 基矢量表示为 e" = . . . :1. . . . ,. e, +. . . :2. . . . ,. e. + . .. :3 R 一 (x\textasciitilde{} + X\textasciitilde{} + x\textasciitilde{}) 川 glf(zf+xi +zi)lmzz+(zf+d+zi 、 i7? e 3 e X 1X3

\[8=(zi+zi)ln(zi+xi +zi)In el\]

十 Z 2 X3 一

\[( x~ + X~) 1/2\]

(x\textasciitilde{} + xD 112 (X\textasciitilde{} + X\textasciitilde{} + xD 112"'2 (X\textasciitilde{} + X\textasciitilde{} + xD 1/2 町 e 一 -Z2+ZI ，一 (x\textasciitilde{} + xD 112"'1 I (xf + xD 1/2"'2 拉梅系数表示为

\[hR = 1,\]

基矢量的导数表示为

\[h8 = R , h", = Rsin()\]

4. JeR E百一叼， JeR ." E石= Sloue" Je8 百一 "R Je. \_ -二= costJe\_

\[O (/J\]

\textasciitilde{}=一 (cos(\}eB + Sin(\}eR)

\[O (/J\]

生旦 =EEE= 生2= 生2=()JR JR JR J(\} - 常用向量导数公式在正交曲线坐标系中的表达式

\[V = e; J- -J = e; -:- ---, hj Jqj\]

附 录

附录 Vw = e; J- ${\rm e}$hp = e, \textasciitilde{}旦旦 +e. \textasciitilde{}旦旦 +e. \textasciitilde{}旦旦rp 一 'hsaqalh18qITZhzaqz TZhaq3

\[VIaa l aa+laa+13a\]

a = e; 瓦元;=el 瓦再;于 e2 瓦再;于 e3 瓦石; 1 ${\rm e}$Ja 1 r ${\rm e}$J ,.. "${\rm e}$J,.. "${\rm e}$J,.. , l = e, ${\rm •}$ 一一- -一一一一|一一 (h 2 仇 al) + 一一 (h ， h 3 a2) + 一一 (h ， h 2 向 ) I , h; ${\rm e}$Jq; h 1

\[hzh3 LeJq, ..-..-.-.. , eJqz ..-....-.. , eJq3 ..-,'.,. J\]

h,e, hzez h3e3 3 a VX a = e; X 一h; ${\rm e}$J q; = h ,hzh3 ${\rm e}$J ql ${\rm e}$J qz ${\rm e}$Jq3 h ,a. hzaz h3a3 I\textasciitilde{}\{ 位主! ${\rm e}$J rp '..J... \textasciitilde{} \{h.h3 ${\rm e}$J rp '..J... \textasciitilde{} (包旦旦旦旦门句=瓦瓦瓦 L ${\rm e}$J\textasciitilde{}， \textbackslash\{\} h, - ${\rm e}$J\textasciitilde{}.)+ ${\rm e}$J $\tilde{z}$ \textbackslash\{\} hz - ${\rm e}$J $\tilde{z}$)+$\theta$$\tilde{J}$ h3 - ${\rm e}$J \textasciitilde{}3) J r L'" , az I L ${\rm e}$J h. L ${\rm e}$J hz \textbackslash\{\} I a3 I L ${\rm e}$Jh , L ${\rm e}$Jh3 \textbackslash\{\} l b ${\rm •}$ Va = I b ${\rm •}$ Va , + 一\textasciitilde{}(b ， 一一一仇一\textasciitilde{} 1+ . ...\textasciitilde{} (b , :"1 - b， 一$\tilde{lle}$, L 且 h 1 h 2 飞 ${\rm •}$ ${\rm e}$Jqz -. ${\rm e}$Jql I ' hlh3 飞 ${\rm •}$ ${\rm e}$Jq3 -. ${\rm e}$J ql I T' + r b. Vaz +主\textasciitilde{}(b2 坐!-b. 坐1. )+主\textasciitilde{}(bz 坐! -b3 旦旦门 zL - . -. , h ， h z 飞 $\theta$ q. -. ${\rm e}$Jq3 I . hz h3 飞 ${\rm e}$J q3 -. ${\rm e}$J qz I T' r L ... $\tilde{I}$ a. I L ${\rm e}$Jh3 L ${\rm e}$J h. \textbackslash\{\}, a3 I L ${\rm e}$Jh3 L ${\rm e}$Jhz \textbackslash\{\} l + I b. Va3 +$\gamma$\textasciitilde{}(b3 了一 -b， 了一 )+T?i b3J-b27i |ez L nln3 飞 c1 q. c1q3 I n2 n3 飞 c1 qz c1 q3 I J 5. 流体力学中常用公式 (1)直角坐标系 (x ， y ， z)

\[U = ui + vj + UIC\]

: ${\rm e}$J rp , : ${\rm e}$J rp , L ${\rm e}$J rp Vrp = i 一 +j 一 +k 一$\theta$ x - ${\rm e}$J y ${\rm e}$J z ${\rm e}$J u ${\rm e}$J v ${\rm e}$J w V. U = 一+一+-=${\rm e}$J x ${\rm e}$Jy ${\rm e}$J z A$\psi$= 匀+匀+匀${\rm e}$J x. ${\rm e}$J y. ${\rm e}$J Z. (${\rm e}$Jw ${\rm e}$J v 飞 (${\rm e}$J u ${\rm e}$JW\textbackslash\{\}.. ,${\rm e}$Jv ${\rm e}$J U\textbackslash\{\} V $\times $ U=132-32t+l 二- ${\rm a}$\textasciitilde{})j + t 瓦 -37) /$\theta$ u ， ${\rm e}$J u , ${\rm e}$J u 、. . I ${\rm e}$J v , ${\rm e}$J v , ${\rm e}$J v 、 U$\cdot $VU=(UEZ+u 巧 + w ;\textasciitilde{} )i + ( u ;; + v 巧

\[+w ;;)i\]

I 习 w ， ${\rm e}$JW. ${\rm e}$J w 、

\begin{gather*}
+ ( u -:\\
+ v -:\\
+ w -:
\end{gather*}

1k 飞 $\theta$ x ' - ${\rm e}$Jy

\begin{gather*}
, - eJ z J\\
\partial ,U = \partial ,ui + \partial ,vj + \partial ,uIc
\end{gather*}

(2) 柱坐标系 (r ， O ， 。

\[u = Urer+USeS+U=e.\]

录 ve=ffe+iEEen 十纽e , Jr' r ${\rm e}$J()' ${\rm e}$J z' 附400 ${\rm e}$JU, I U, I 1 JUo I JU. V. U = 一二十:::....-:+一..::\_. +一一 ${\rm e}$J r . r . r J() . Jz I 1 au. auo 飞 I IJU , aU. \textbackslash\{\} VXU=( 一一一 一::..!!.le. + (一一一一:-. le. 飞 r J() r1 z r' . 飞 az ar r" ab 匹柏l-r u7 + 叫-h+ -N + 世。一柑1-2r + J-h I-r + -N 一A 2-t-Z J-3 + 坐柑+ 组

\[h+ ZA--r 2u\]

一、。。，A

\[l.rT U, 2 aun  \]

ll.U = 1 ll. U， 一$\rightarrow$一寸$\leftarrow$=-\_. 1 飞 r" r' c1() I ，中'U A + ，中' 叫一川队

\[-M\]

叫 气d-a2-zr + U A ,,, IrT au, I Uo JU, , rT iJU, U\textasciitilde{} \textbackslash\{\} U.VU=IU 一二+\textasciitilde{}.一二 +U. \textasciitilde{}一一\textasciitilde{}. le. 飞 iJ r . r J() ${\rm ‘}$ r1 z r r' 1. T iJU. . U. iJU. . ${\rm ••}$ JU. . U.U. 飞+IU，"::'+::::.\_<! 一:-\_" + u. ..::" + \textasciitilde{} ，\textasciitilde{}. le. 飞 ， Jr . r iJ() . -. c1 z' r I $\sigma$ I'T JU\_ . U. aU\_ . 'T ${\rm e}$1U. \textbackslash\{\} +IU 一二十」$\rightarrow$:-\_' +U 一$\tilde{l}$ 飞 Jr . r a() . -. az I (3) 球坐标系恨，()， $\psi$)

\[U = UReR + Uoeo + Ufe~\]

\_ JI${\rm þ}$ \_ I 1 r11${\rm þ}$\_ I 1 al${\rm þ}$ W 一 \_ :, . ...e R 十一 \_T\_.e O + 一一一 -'-eJR-" , R a()-' ' RsinO a\textasciitilde{}口伊 iJUR I 1 JUo I 1 JU, I 2UR I Uocot() V. U = 一一+一一一+一-一--， +一一+一一一一JR ' R J() , Rsin() iJ cp 'R R f 1 a 口 l 坠旦!:!!\_ \_ \_1\_ auo\textbackslash\{\} VXU =(一一\textasciitilde{}+ 一一一一一::..!!. le \textbackslash\{\} R iJ() I R Rsin() a伊/町 (Rs\textasciitilde{} 机 a u.\_ U 、/肌 Uo 1 auR\textbackslash\{\} 一一一一」一:..::::..!f一=-r! )eo + I 一:: + :: 一一$\tilde{n}$" ) Rsin() a1> JR R r. ' \textbackslash\{\} JR 'R R r10 r'" a2 I 2 a I cot() J , 1 a2 I 1 a2 ll.=一一+一一+一一一十一一一+一一一一一一$\theta$ R 2 ' R JR ' R 2 ${\rm e}$J() ,. R2 J()2 ' R 2sin20acp2 A$\psi$= 勾+主业 + c;>\_\textasciitilde{}() \textasciitilde{}业+土当+$\leftarrow$上 a 2 业 r1R2 . R JR . R 2 r1() , R2 r1\textasciitilde{} , R2sin20 Jcp2

附录 l.rT 2UR 2 ${\rm e}$J( UosinO) 2 au" \textbackslash\{\} ${\rm ð}$.U = I ${\rm ð}$.UR 一一一一一一一-一--一一一一一一一一一一 Ih

\[  ~\tilde{K} R R 2sinO ao R2sinO a \psi  r\]

K /.TT I 2 aUR Uo 2cosO au",\textbackslash\{\} + I ${\rm ð}$.U. + 一一一一一一一一一一一一一一::....! le. \textbackslash\{\} \textasciitilde{}\textasciitilde{}. , R2 iJO R2sin 20 R2sin 20 a $\psi$/ 。 \{.rT I 2 rJ的 2cosO r1Uo 一旦f..\_\_\_ \textbackslash\{\} + I ${\rm ð}$.U m + 一一一一一+一一一一一一-\_.，-.'.， \_l飞" ' R2sinO aq; 'R2sin 20 r1q; R 2sin20r" ITT r1UR , U. auR , Um r1UR U\textasciitilde{} + U! 、U.VU=IUR 一\textasciitilde{}.: + :"一\textasciitilde{}\_K +一:::..f\_ \textasciitilde{}一 $\tilde{Ih}$ 飞 " aR ' R ao 'RsinO r1 $\psi$ r l< ITT aU. , U. au. , Um aU. , URU. U! 俨 otO 飞+ I U R \textasciitilde{}\textasciitilde{}\textasciitilde{}..: + :'"一\textasciitilde{}-" +一=-..f..\_一\textasciitilde{}" + \textasciitilde{}':.\textasciitilde{}"一::::....I!.:一$\tilde{I}$飞 R JR R JORsinOJ$\psi$ R--R-j 1.. iJU\_ , U. au\_ , U\_ au\_ , U\_U" , U.U 飞+ I UR .. \textasciitilde{}-\textasciitilde{}， +:" ..-:\_, +一'::..!---;.\textasciitilde{}:..+\textasciitilde{}$\tilde{K}$+::::\_$\iota$$\tilde{cot}$Ole 飞 n aR ' R JO 'RsinO a\textasciitilde{}口 R ' R \_\_\_V I N 流体的变形速率张量和牛顿流体本构方程 流体的变形速率张量和应力分别用 s= \{Sij \} .p= 怡ij\}(i.j=1.2.3) 表示，它们 之间的关系称为流体的本构方程.牛顿流体的本构关系式为 $\rho$。= [- p+ ( $\mu$' 一 3钊$\mu$片)su矗H牛矗 式中 $\rho$ 为压强， $\mu$ 是剪切粘性系数 '$\mu$ ，为体积粘性系数，又称第二粘性系数. 下面给出一般曲线坐标系和常用坐标系中以上关系的表达形式. 1.正交曲线坐标系(钮 ${\rm •}$ q2 ${\rm •}$ q:l) 中变形速率张量和本构关系的表达式 令速度场为 : U=U1el +U2 e 2+U:， 酌，变形速率张量各分量如下 z

\[1 aU1 I U2 \theta  h l I U3 ah 1\]

.$\tilde{II}$ =瓦石7T 瓦瓦石;于瓦瓦石; h=1 但主+卫L 壁旦+卫L 壁1.'22 - h2 r1 q2 I h2 h3 rJq:. I h1h2 aql 1 r1U3 I U1 ah:, I U2 ah:, .\textasciitilde{}:，:， =瓦石;$\uparrow$瓦瓦石:于瓦瓦药;

\[1 rh z a ,U2   I hl iJ U 1   l\]

.\textasciitilde{}12 = S21 - 2L$\tilde{J}$百飞瓦 f 瓦药;\textbackslash\{\}瓦 )J 3 = .\textasciitilde{}:!2 = \textasciitilde{} [主主(去)+去元(去 )J :, = s:!1 - \textasciitilde{} [去未(号)+去看(去 )J

牛顿流体的本构关系如下 z >$\rho$扣11 = 一 $\rho$ +2严 1川1 + ($\leftarrow$$\mu$ , 一 3卡妇$\mu$片)户〈 句川+5句挝22 + s句阳M3:1川l) $\rho$扣22 = 一 $\rho$ +2严阶内22什+($\leftarrow$$\mu$, 一 3争妇$\mu$片)户( 归川+ 句扭川2+ 如 $\rho$ 3臼3= 一 $\rho$ +2 严阶3计 ($\leftarrow$$\mu$ , 一 $\div$轩妇$\mu$片)严(5仙问51归1川1+ 川 533臼3 ) PI2 = P21 = 2ps l2 '$\rho$23 = P32 = 2ps 23 , P31 = P13 = 2四" Z 直角坐标系中的变形速率张量和本构关系 令速度场为 :U=U.r e .r +Uyey+V.e. ， 变形速率张量各分量如下 z aU. aU.. aU .rx = a:r' Syy = a歹 ， S.. = aZ二

\[1 ,aU:r I au. \]

Szy = 5,.. =言 $\tilde{a}$歹 '33fj 1 ,au. I au. 飞 5 )0' = S\%JI = 言飞 aZ千-r- ay )

\[1 ,au. I au.r  \]

\textasciitilde{}$\tilde{a}$ = S.r.z = 王飞 3言-r a; ) 直角坐标系中牛顿流体的本构关系如下 z P:r:r = 一 $\rho$+2灿 +($\mu$'-3$\mu$) (s.r.r +Syy +S..) pyy =- p+2内 +($\mu$'-3$\mu$) (s.r.r +Syy +5..) 户自 =-p+2内 +($\mu$'-3$\mu$) (S.u + Sy.y + 5.. )

\[P.ry = \rho yr = 2ps.ry' Pyo = \rho '.V = 2psyo , Pu = \rho n = 2psu\]

3. 柱坐标系中的变形速率张量和本构关系 令速度场为 : U=U.e.+U${\rm ø}$e${\rm ø}$+ U儿，变形速率张量各分量如下 z S.. = 笠E a-lau t 旦 S.. = 空三 \textasciitilde{}'咱 -73百 I r' '.. - az S.${\rm ø}$ = \textasciitilde{} [手夺 +r (于 )J S${\rm ø}$. = \textasciitilde{} [r 乏(子)+$\downarrow$等] S.. = \textasciitilde{} (等+尝) 附录

附录 柱坐标系中牛顿流体的本构关系 $\rho$ rr - 一 $\rho$+2四 "+($\mu$'-3$\mu$ )(S" +SI16 十 S.. ) $\rho$116 =- p+2内 +($\mu$' 一切 (S" + SI16 + S..) $\rho$..- 一 $\rho$+2$\mu$..+($\mu$'-3$\mu$ )(S" +SIlI +S..) $\rho$ rlI = $\rho$如 =2$\mu$5 rl1' PO. = $\rho$ =0 = 2$\mu$50. , P.. = P.. = 2严旷 4. 球坐标系中的变形速率张量和本构关系 令速度场为 : U=UReR+UoeO+ U"，鸟，变形速率张量各分量如下 z aUR 1 auo I UR SRR = aR' 5116 = R a百 TE 1 aum , UR , U.CO舍A =一一一 - \_:-， + -\textasciitilde{}" + -.;一一RsinO a(/J , R' R r $\tilde{a}$ ,U. \textbackslash\{\}, 1 aUR l SRO - 一 1 R \_\textasciitilde{}$\tilde{r}$ :: 1+ 一 \textasciitilde{}:"lC I 2 L-- aR 飞 R l' R ao J 1 rsinO a I Ue 飞 I 1 auo 1 50", = zL R aO\textasciitilde{} 击。)-r Rs币。可」

\[1 r 1 aU R , \tilde{a} ,U m   l\]

R =一|一一一 -::K +R \_\textasciitilde{}$\tilde{r}$\textasciitilde{} 11 2 LRsinO a(/J ,-- aR 飞 R IJ 球坐标系中牛顿流体的本构关系 PRR =-$\rho$+2阳 +(i-3$\mu$ )(SRR +s押 +$\omega$ $\rho$押=一 $\rho$+2内 +($\mu$'-fd(sm+sw+ 韧〉 阳=一 $\rho$+2内 +($\mu$' 一切 (SRR +s""" + 抽) $\rho$R" -$\rho$州 =2同 R" '$\rho$\textasciitilde{} = Po" = 2间\textasciitilde{}， POR = $\rho$ RB = 2PSOR V 牛顿型流体运动的基本方程(纳维，斯托克斯方程) 1.向量形式的基本方程 连续方程

\[an ... ~ + v . (pU) = 0 at\]

404 附 运动方程 iJU 一 +U.VU=f+$\tilde{V}$.piJl ,- - - '" p 能量方程 I I U 2

\[  1.' ..., I I U\]

2 \textbackslash\{\} ,..., I 1 ..., ,.... ..,.... 1. I ${\rm A}$ \textasciitilde{}\_Ie+\textasciitilde{}.... I+U.VI f'+\textasciitilde{}.... l=f.U+$\tilde{V}$.(p.U)+$\tilde{q}$'+\textasciitilde{}${\rm ð}$T \textbackslash\{\} - , 2 I 飞 2 I J -, p' ,- -" p" p 状态方程

\[\rho  = f(p. T)\]

式中 f 为单位质量的质量力 .p 为密度 .U. VU 为迁移加速度 .V ${\rm •}$ p 为应力张量的 散度 .q 为单位体积的热源. 2. 应力张量散度的表达式 (1)一般曲线坐标系中应力张量的散度表达式 r 1 r iJ".." iJ".." J".. , l V;P=\textasciitilde{}-. -:-.1 一一 (h 2 h 3 $\rho$ 11) +一一 (h 3 h l $\rho$ 12) +一一 (1I 1 h 2 $\rho$ I:i) 1 \textbackslash\{\}h lh 2h" LrJql '--.--,'..-", , c1q2 "-,'--'..-'.' , Jq;! "-, -...-'," J 1 Jh 1 1. 1 Jh 1 1 Jh 2 1 Jh 几 i+$\rho$l 一一\textasciitilde{}+如一一一 - P22 一一一 -P:u 一一 $\tilde{t}$1 :i $\tilde{zlt}$lh22qzIh1123Jq3hlhzJql hlhBJqlJ +一[(h$\iota$川2(1J J

\[h lh 2h3 LJql "-.--,'..-"- , Jq2 ---0--' ,,-..- , r1q;! t\]

1 Jh z 1. 1 r1h 2 1 rJ h" 1 Jh l \textbackslash\{\} +$\rho$2;! 一一 1\textasciitilde{}"2 +$\rho$I 一一一一一户"一一一一一 $\rho$1 一一一:，qh2 h 3 $\theta$q32hzltl3qIhzhzaqzl hzltl3qzj

\[(I J J + h ,  h\]

1 .IIJ 一仙 2 h 3P;!I) + \textasciitilde{}(川2;!) +立叫川]hlh2113 LrJql '--.--""-'" - , iJqz '--,'--' ..-." - , rJ q] 1 c1I1,! 1. 1 r1h ;! 1 ${\rm e}$Jh 1 ${\rm •}$ 1 JII2 \textbackslash\{\} +如一\textasciitilde{} \textasciitilde{}".i +$\rho$ 一一一 -PII 一一一一 $\rho$ 一一 '\textasciitilde{}"2 ) 1 1I ;!hJ ${\rm e}$J ql 剖 112 h 3 $\theta$ q2 1' 11 hlh3 Jq;! 1'22 11 2 lI;j dq;j f (2) 直角坐标系中应力张量散度表达式 r CJ ${\rm þ}$ ,r.r I d ${\rm þ}$ ,r ,' I CJ ${\rm þ}$ .r, 1 V. P=I 二王+二2:. +$\iota$一 le.

\[L J I ' rJ Y rJ z J-'\]

rdP ,r.v I cl${\rm Þ}$y.v I d${\rm Þ}$.r. l\_ I rcJ ${\rm þ}$ ", I cJ ${\rm þ}$". I d${\rm þ}$ ", 1 +1 亏-，r.v +气且+ "\textasciitilde{}.r. le v + 1 七三+气'". +二$\tilde{Z}$\% I

\[L (/I (/ Y (/ Z J - L (/.. (/ Y (/ Z J\]

(3) 柱坐标系中应力张量散度的表达式 \{cJ ${\rm þ}$" I 1 r1 p.o I rJ p$\rho$w 一 $\rho$00 \textbackslash\{\} V. P=I 一工+一一一旦+一三+一一一一'f!I!. le. 飞rJ r ' r JO 'iJ z r r' /$\theta$$\rho$ .0 , 1 rJ ${\rm Þ}$oo I J Po. I 2$\rho$ .0 \textbackslash\{\} +1 一一+一一\textasciitilde{}$\tilde{o}$ + 一\textasciitilde{}(/\textasciitilde{} +一'J rO le n 飞 dr ' r c10 ' rJ z ' r J $\sigma$ \{J P.. , 1 ${\rm e}$J ${\rm Þ}$o. I d P $\rho$.. \textbackslash\{\} +1 一二+一 v 1' 0三+一二三+一'.':.le 飞 cJ r r cJ(\} dZ r I 录

附录 405 (4) 球坐标系中应力张量散度表达式 V. p =1 旦旦旦+ :. \textasciitilde{}担+ \textasciitilde{} \textasciitilde{} \_ J $\tilde{qII}$ + :. (2 阳 - ${\rm þ}$lB - ${\rm þ}$'fXF + ${\rm þ}$ llfJ cot() 门 eRL iJR ' R iJ() 'Rsin() J fP ' R '-1'''''' yw Y 'fXF' ynu--- v ' J r J ${\rm þ}$ Rli I 1 ${\rm e}$J ${\rm þ}$ lB I 1 J ${\rm Þ}$o", I 1 I \textasciitilde{} \textasciitilde{}'\_\_.n I 3 $\tilde{l}$ +1 一一十一 U.:\textasciitilde{} + 一一一-」 +-Wm 一 $\rho$ 'fXF )cot() + \textasciitilde{}${\rm Þ}$ RIi leo L JR ' R J() , Rsin() J fP ' R 'yw Y 'fXF '---v , R Y "" J +1 旦旦旦旦+ :. J\textasciitilde{}$\tilde{f}$ + \textasciitilde{} ${\rm e}$J\textasciitilde{}'f'f + :. (忡'f R + 2${\rm þ}$o'l' co 叫eL iJ R ' R J() 'Rsin() J fP ' R '-Y'f" ' -ro'l' ---V' J 3. 不可压缩牛顿型流体运动的基本方程 (1)向量形式的基本方程 连续方程 V ${\rm •}$ U = 0 运动方程 JU ，..\textasciitilde{}. V b $\tau$ +U. vu =一二.x + f+ II!:lU rJ[ P 能量方程 字 +U-vr=i v.(AVT)+$\phi$ +cj 01 P 式中 $\phi$ 是耗散函数. (2) 直角坐标系中的基本方程 连续方程 nu --拟二、"+ 也句+ 以-h 运动方程

\[JU, , rT JU" , rT au" , rT JU , 1 J\rho \]

一\textasciitilde{}' +U 一」十 U 一\textasciitilde{}' +U 一一=一一$\rightarrow$+!， +$\nu$!:l U" al ,-, J.r '-.v Jy ,- Z Jz P J.r JU" , ${\rm ••}$ aU" , ${\rm ••}$ JU" , ${\rm •}$ T aU" 1 J b 一:，v +U 一-.v +U 一-\_v +U 一::::...z.\_一\textasciitilde{} $\tilde{P}$ + Iv +$\nu$ !:l U，. rJ l '-," J.r '-,V rJy , -. .Jz P Jy aU \_ . ${\rm ••}$ JU \_ , ${\rm •}$ T aU \_ , T' JU $\epsilon$1 J${\rm þ}$ 一二 +U 一:::...:.+U 一」十 U 一一 +!. +$\nu$!:l U\_ .11 ,-, cJ.r ' -." rJ y ,-. rJ z p .J z 能量方程 Je , r r Je )" 一十 U$\tau$一 +U ， $\tau$一 +U， \textasciitilde{}. = \textasciitilde{}!:l T+ $\phi$ +d rJ l ' -. .1 x 、 cJ y --az p ${\rm •}$ 式 JIJ $\phi$=$\mu$\{2[( 等 f + (芳$\tilde{f}$ + (安门 +$\mu$[( 尝+毛主 f 十(去+尝 f + (弩+芳$\tilde{f}$ J\}

式中 (3) 柱坐标系中不可压缩牛顿型流体运动的基本方程 连续方程 运动方程 能量方程 arU. , au. , aU 一一工+一-旦 +r 一\textasciitilde{} = 0 ar ' aO ' . (\}z 2 盯 J 盯盯o au, I r T aU, U\textasciitilde{} 一\textasciitilde{}' +U， 一:::..!.+:::..!!一-\_' + U. 一一一一

\[\theta  t . -. ar . r ao ' (z r\]

1 a ${\rm Þ}$ , r , 1. r T U, 2 auo \textbackslash\{\} =一-=- ':1' + f , + $\nu$ (${\rm ß}$U 一一一一 -::"U ) p ar 飞 ， r 2 r aO I aU. , TT aU. , U. aU. , TT au. , U.U. ..::. + u 一::..!+:::..!!一-\_. + u. 一::..!+一\textasciitilde{}

\[at . -. ar . r ao . - < az' r\]

1 a b , r , 1... Uo , 2 aU, \textbackslash\{\} =一一二$\iota$ +fo+ $\nu$ (${\rm ß}$Uo 一 -2' +$\tau$ 一-l p rao 飞 r 一 r' ao I aU. I TT au. I Uo au. I rT au. 一\textasciitilde{}. +u， 一\textasciitilde{}. +:::..!!一\textasciitilde{}\_' + u. 一-at . -. ar . r ao . - < az 1 ab =一-=- "\_, y + f. + $\nu$ ${\rm ß}$u.

\[\rho  <1 Z\]

ae , T T ae , U. ae $\theta$ e ). 一 +u 一+:::..!!一 +u 一 =\textasciitilde{}${\rm ß}$T+ $\phi$ +cjat ' -, ar ' r ao ' -. (\}z p $\phi$ = 2p.[ ( 安 f + (;芳+号 f + (安门 附录 十 $\mu$[( $\downarrow$等+尝 f + (尝+等 )2+(fZE+ 号子一千 fJ (4) 球坐标系中不可压缩牛顿型流体运动的基本方程 连续方程 运动方程 1 a , \textasciitilde{}. T T " 1 a'T T ${\rm •}$ \_" 1 I au m \textbackslash\{\} 寸一 (R'UR ) + 一一一 (UosinO) + 一一一$\leftarrow$:1. 1=0R2 aR'-' -". , RsinO aO.-...... v ${\rm •}$ , RsinO 飞 a\textasciitilde{}口/ 川 川口 3 门卫LJUR Uj土豆;U$\tilde{R}$ +UR "\_:-'\_! + ':.,0 U\textasciitilde{}: + 一一一

\[\theta  t ,-" aR ' R ao 'R sinO ehp\]

1 a${\rm Þ}$ I r I r .rr 2UR 2 a(Uo sinO) 2 缸里丁=一 -EZ+h+v|AUR 一一手 - n2.\_'\_一一一一一一一一 | p aR I JK I "L

\begin{gather*}
~-K\\
R 2 R2sinO ao R2sinO a(/J J
\end{gather*}

aU\_ . TT aU\_ , U. aU\_. U\_ aU\_ , U的rT 1T< T 一\textasciitilde{}+UR 一\textasciitilde{}+ -:\_.一=="+一\textasciitilde{} --::f + -:-, + -\textasciitilde{}'cotO

\[at ,-" aR ' R ao 'RsinO a(/J' R R\]

1 1 a ${\rm Þ}$ , r , r. T T I 2 aU R U m I 2cosO aU门 =一一一\textasciitilde{}\_l] $\tilde{p}$ + f ., +111 ${\rm ß}$U., + 寸一」一---=....!..一+一一一|p RsinO a $\psi$ .， I "L

\begin{gather*}
~-.,  I\\
R2sinO a(/J R2sin 20
\end{gather*}

I R2sin 20 a 伊」

附录 式中 ${\rm e}$JU, I rT ${\rm e}$JU8 I U, ${\rm e}$JU, I U" ${\rm e}$JU${\rm ø}$ I URU, U! 一:'， +UR 一一 + V\_\_' "::\_' +一tJ! \_ "\_:: 8 + ""'!\_'-' B 一${\rm e}$J t ,-'" ${\rm e}$JR ' R ${\rm e}$JO 'Rsi nO ${\rm e}$J rp' R R 11 ${\rm e}$J ${\rm Þ}$ , r , f. rT , 2 ${\rm e}$JUR Ue 2co sO旦旦丁 =一一:， ':.\textasciitilde{} + f , + $\nu$I ${\rm ð}$.U, + 寸一一一一一一一-一一 | p R ${\rm e}$JO I JU I '1 ...\textasciitilde{}. , R2 ${\rm e}$JO R2sin20 R Z sin20 ${\rm e}$J $\psi$ 」 能量方程

\[eJe 'FT eJe , Un eJe , Um eJ e \lambda  , \]

一 +UR 一 + -:u 一+一::...f.\_一=.!!.. V' T+ $\phi$+4 ${\rm e}$J t ' -'" ${\rm e}$JR ' R ${\rm e}$JO ' RsinO ${\rm e}$J rp p $\phi$=$\mu$\{2[ (尝 f + (古等+生 f + (对函弩+安+哈望 fJ +[击。专+乎是 (\textasciitilde{})J +[击。写 +R ${\rm e}$J\textasciitilde{}( 去 )J + [R 录(导)+古等 J\}

d 附 表 1.常见液体物性参数表

\[\rho = 1. 0132 X 105 Pa.t= 20"C\]

p f.l 液体

\[/<kg' m- 3) /<l O'kg'm-l.s 1)\]

水 998 10. 1 苯 895 6.5 四氯化碳 1588 9.7 汽油 678 2.9 甘油 1258 14900 液氢 72 0.21 煤油 808 19.2 水银 13550 15.6 液氧 1206 2.8 SAEI0 918 820 SAE30 918 4400 Y 1阳 、 E /<N'm- I ) /10 :1 Pa /OO-'N'm 勺 0.073 2.34 2070 0.029 10.0 1030 0.026 12.1 1100 55 0.063 14 X 10 ${\rm •}$ 4350 0.003 21. 4 0.025 3.20 0.51 17X 10 26200 0.015 21. 4
